\chapter{Cloud-Native and Kubernetes Storage}
\label{ch:15}

The containerization movement, led by Docker and Kubernetes, redefined not only how applications are packaged and deployed but also how storage is consumed and managed. Classical applications ran on VMs or physical servers with persistent filesystems that were assumed to outlive any individual process; containers, by contrast, are ephemeral by design --- a container's filesystem is discarded when the container exits, and a new container spawned from the same image starts with a clean slate. This ephemerality is a feature for stateless applications (it enables reliable, reproducible deployments), but it creates a fundamental tension for stateful workloads --- databases, message queues, analytics engines --- that need storage that outlives the container lifecycle.

This chapter examines how cloud providers have built block and file storage services that integrate with compute infrastructure at platform scale, how Kubernetes has developed a sophisticated abstraction layer for storage that decouples application definitions from storage implementation, how container-native storage vendors have built distributed storage systems specifically for Kubernetes environments, and how multi-cloud storage strategies can provide data portability across cloud boundaries.

\section{Cloud Block Storage}
Cloud block storage provides raw storage volumes that attach to virtual machine instances, functioning analogously to a local disk from the operating system's perspective. The operating system sees a block device, formats it with a filesystem, and issues reads and writes that the cloud provider translates into operations on its distributed storage infrastructure. The storage itself is not on the VM's host; it resides in the cloud provider's storage backend (typically large-scale distributed systems using thousands of commodity drives with erasure coding) and is presented to the VM over a high-speed internal network.

\section{Amazon Elastic Block Store (EBS)}
Amazon EBS is the foundational block storage service for AWS workloads. EBS volumes are provisioned independently of EC2 instances and attached over AWS's dedicated storage network (the same physical network infrastructure but logically separated from instance-to-instance traffic). The current EBS architecture uses a replicated storage backend with synchronous replication within an Availability Zone, providing durability of 99.999\% (five nines) or better.

EBS offers several volume types targeting different performance profiles:

\textbf{io2 Block Express} is the highest-performance tier, delivering up to 256,000 IOPS and 4,000 MB/s throughput per volume, with sub-millisecond latency and 99.999\% durability (compared to the standard EBS 99.8-99.9\% for older volume types). Block Express uses NVMe hardware and supports IOPS provisioning independently of volume size, allowing workloads to specify exactly the performance they need. io2 Block Express supports \textbf{Multi-Attach}: a single volume can be simultaneously attached to up to 16 EC2 instances in the same Availability Zone, enabling shared block access for cluster applications (clustered databases like Oracle RAC, for example) that require shared storage.

\textbf{gp3 (General Purpose SSD)} is the standard choice for most workloads, providing up to 16,000 IOPS and 1,000 MB/s throughput with the ability to independently provision IOPS and throughput without increasing volume size. gp3 replaced gp2 as the recommended default, eliminating the gp2 model where IOPS scaled rigidly with volume size (3 IOPS/GB with burst).

\textbf{st1 and sc1} are HDD-based volume types optimized for throughput-oriented workloads (Kafka, log processing, large sequential scans) and cold storage respectively, offering high throughput at lower cost than SSD tiers.

A critical architectural characteristic of EBS is that volumes are \textbf{Availability Zone-local} --- a volume created in us-east-1a cannot be directly mounted on an instance in us-east-1b. Cross-AZ access requires creating an EBS snapshot (stored in S3 and therefore region-wide), creating a new volume from that snapshot in the target AZ, and attaching the new volume. This AZ locality simplifies the consistency model (all replicas are within the same physical data center facility) but requires attention in multi-AZ application architectures.

EBS snapshots are incremental and stored in S3. Each snapshot captures only the blocks that have changed since the previous snapshot, making the first snapshot full and subsequent snapshots much smaller. The AWS Recycle Bin and EBS Snapshot Lifecycle policies allow automated snapshot management, retention, and cross-region copying for disaster recovery.

\textbf{EBS-optimized instances} are EC2 instance types that provide a dedicated network path between the instance and EBS storage, ensuring that EBS I/O does not compete with instance-to-instance network traffic. Nearly all modern EC2 instance types are EBS-optimized by default.

\section{Azure Managed Disks}
Azure Managed Disks are the Azure equivalent of EBS, providing persistent block storage for Azure VMs. Before Managed Disks, Azure used "unmanaged disks" --- page blobs in Azure Storage --- which placed the burden of storage account management (including IOPS limits per storage account) on the administrator. Managed Disks abstract this away: Azure manages the storage account placement, replication, and capacity, and users interact only with the disk object.

Azure Managed Disks offer four performance tiers:

\textbf{Ultra Disk} is Azure's highest-performance offering, providing up to 160,000 IOPS and 4,000 MB/s throughput per disk, with configurable IOPS and throughput independently of disk size. Ultra Disk latency is sub-millisecond and it supports dynamic performance adjustment without downtime --- IOPS and throughput can be increased or decreased through the Azure API without detaching or resizing the disk. Ultra Disk is the appropriate choice for tier-1 databases (SQL Server, Oracle, SAP HANA) with predictable I/O demands.

\textbf{Premium SSD v2} extends Premium SSD's IOPS-per-GB model with independent IOPS and throughput configuration (similar to Ultra Disk), sub-millisecond latency, and broader regional availability compared to Ultra Disk. Premium SSD v2 is generally the preferred choice for most database and production application workloads.

\textbf{Premium SSD} (v1) is the original SSD tier, delivering up to 20,000 IOPS and 900 MB/s per disk, with IOPS scaling with disk size (not independently configurable). Premium SSD is appropriate for medium-performance workloads.

\textbf{Standard SSD and Standard HDD} serve lower-performance workloads where cost efficiency matters more than latency.

Azure Managed Disks support \textbf{shared disks}: a single disk can be attached to multiple VMs simultaneously (up to 10 VMs for Premium SSD and 5 for Ultra Disk), enabling shared storage for Windows Server Failover Cluster (WSFC) deployments, SAP ASCS clustering, and Oracle RAC. The application using shared disks must implement its own distributed locking (SCSI Persistent Reservations are supported) since the disk itself provides no coordination.

\textbf{Azure Disk Bursting} allows disks and VMs to burst beyond their provisioned IOPS and throughput limits for short periods, accommodating periodic spikes in I/O demand without requiring over-provisioning.

\section{Google Cloud Persistent Disks and Hyperdisk}
Google Cloud's block storage evolution has moved from Persistent Disk (PD) to the more capable \textbf{Hyperdisk} architecture.

\textbf{Persistent Disks} are the established block storage service, offering zonal and regional variants. Zonal PDs are replicated within a single zone; Regional PDs are synchronously replicated across two zones in the same region, providing higher availability for stateful VMs at the cost of write latency (a write must be acknowledged by both zones). PDs support snapshots (stored in Cloud Storage), live resizing, and multi-reader attachment (a single PD can be attached to multiple VMs in read-only mode simultaneously).

\textbf{Google Cloud Hyperdisk} is the next-generation block storage platform, designed with independently configurable throughput, IOPS, and capacity. Hyperdisk delivers up to 350,000 IOPS and 5,000 MB/s per volume (for Hyperdisk Extreme), with sub-millisecond latency. Hyperdisk volumes support \textbf{dynamic performance scaling} --- changing IOPS or throughput without creating a new disk or detaching --- and \textbf{tiered storage} within the service.

The performance comparison across providers is workload-dependent. {[}Hyperdisk Extreme can deliver up to 350,000 IOPS for read operations{]}(https://cloud.google.com/compute/docs/disks\#comparison), while EBS io2 Block Express tops out at 256,000 IOPS and Azure Ultra Disk at 160,000 IOPS per volume. Real-world performance depends heavily on instance type, networking configuration, and whether the workload requires random or sequential I/O.

\section{Cloud File Storage}
Cloud file storage services provide managed NFS or SMB file shares that can be mounted across multiple VM instances simultaneously, eliminating the single-attachment constraint of block storage. These services are essential for applications requiring shared filesystem access --- content management systems, home directories, lift-and-shift NAS migrations, and stateful containers that need shared persistent storage.

\section{Amazon EFS and FSx}
\textbf{Amazon Elastic File System (EFS)} provides managed NFS (v4.0 and v4.1) with elastic capacity --- users do not provision a size; EFS grows and shrinks automatically as files are added and removed. EFS storage is distributed across multiple Availability Zones within a region, making it resilient to AZ failures. EFS supports two performance modes: \textbf{General Purpose} (lower latency, suitable for most workloads) and \textbf{Max I/O} (higher aggregate throughput for highly parallel workloads at the cost of higher per-operation latency). EFS Intelligent-Tiering automatically moves files that have not been accessed for a configured period to a lower-cost infrequent access storage class.

EFS is appropriate for content repositories, WordPress media, CI/CD shared artifacts, and container applications requiring shared persistent storage. Its elastic nature and AZ resilience make it operationally attractive. The per-GB-month cost is higher than S3, and performance at high concurrency can be a constraint for demanding workloads.

\textbf{Amazon FSx} is the umbrella for AWS's managed file storage services built on proven commercial and open-source filesystem implementations:

\begin{itemize}
  \item • \textbf{FSx for Windows File Server}: A fully managed SMB file server built on Windows Server, supporting Active Directory authentication, DFS namespaces, shadow copies, and the full SMB feature set including SMB Direct (RDMA). It is the natural choice for lift-and-shift Windows applications requiring SMB shares.
  \item \textbf{FSx for Lustre}: A managed Lustre filesystem tightly integrated with S3 as the data lake backend. Files in an S3 bucket can be linked to an FSx for Lustre filesystem, appearing as local files that are lazily loaded from S3 on first access. Subsequent accesses are served from the Lustre filesystem directly. This integration makes FSx for Lustre ideal for ML training jobs that need high-bandwidth parallel filesystem access to training data stored in S3.
  \item \textbf{FSx for NetApp ONTAP}: A managed instance of NetApp ONTAP running in AWS, providing NFS, SMB, and iSCSI access with full ONTAP feature set (snapshots, SnapMirror replication, deduplication, compression, thin provisioning). This is the preferred landing zone for organizations with deep ONTAP expertise wanting to retain familiar data management workflows in AWS.
  \item \textbf{FSx for OpenZFS}: A managed ZFS-based filesystem providing NFS access with ZFS features: instant snapshots, clones, data compression, and predictable low latency.
\end{itemize}

\section{Azure Files and Filestore}
\textbf{Azure Files} provides fully managed cloud file shares supporting both SMB (2.1, 3.0, 3.1.1) and NFS 4.1 protocols. Azure Files shares exist within Azure Storage accounts and can be mounted simultaneously by multiple VMs (or on-premises servers through Azure File Sync). Like Azure Managed Disks, Azure Files offers multiple performance tiers:

\begin{itemize}
  \item • \textbf{Transaction optimized} (HDD-backed): Suitable for workloads with frequent small I/O operations.
  \item \textbf{Hot and Cool}: SMB General Purpose v2 tiers for standard and infrequently accessed shares.
  \item \textbf{Premium} (SSD-backed): Low-latency shares for performance-sensitive applications. Premium File Shares provision capacity up front (unlike standard shares which are pay-per-GB consumed).
\end{itemize}

\textbf{Azure File Sync} is a hybrid synchronization service that allows on-premises Windows File Servers to synchronize with Azure Files, with cloud tiering for files not accessed locally. This provides a migration path from on-premises NAS to cloud file storage while maintaining local access performance.

\textbf{Google Cloud Filestore} provides managed NFS file shares for GCE VMs and GKE clusters. Filestore is available in Basic (up to 500 MB/s) and High Scale (up to 25 GB/s) tiers, with the High Scale tier providing petabyte-scale NFS for large parallel workloads. Filestore Multishares allows multiple NFS shares from a single Filestore instance, improving cost efficiency for multi-tenant deployments.

\section{Kubernetes Storage Architecture}
Kubernetes was initially designed with stateless applications as the primary target --- its scheduling, scaling, and update models are optimized for pods that can be freely killed and recreated. The storage abstraction layer was retrofitted as stateful workloads increasingly demanded Kubernetes orchestration, resulting in a layered model that provides significant flexibility at the cost of conceptual complexity.

\section{Volumes, Persistent Volumes, and Persistent Volume Claims}
In Kubernetes, a \textbf{Volume} is an abstraction over a directory accessible to containers in a Pod. The most primitive volumes (emptyDir, hostPath, configMap, secret) are either ephemeral or tied to the Pod's lifecycle --- they do not survive pod deletion. Persistent storage requires the PersistentVolume subsystem.

A \textbf{PersistentVolume (PV)} is a cluster-scoped resource representing a piece of storage provisioned in the cluster. A PV exists independently of any Pod; it has its own lifecycle managed by Kubernetes, persisting until explicitly deleted. A PV is bound to a specific storage implementation: an NFS share, an EBS volume, a Ceph RBD image, a CSI-provisioned resource, and so on. PVs have an \textbf{access mode} that specifies how the storage can be mounted:

\begin{itemize}
  \item • \textbf{ReadWriteOnce (RWO)}: The volume can be mounted as read-write by exactly one node. This is the most common mode, corresponding to typical block storage semantics.
  \item \textbf{ReadOnlyMany (ROX)}: The volume can be mounted as read-only by many nodes simultaneously. Useful for distributing static data to many pods.
  \item \textbf{ReadWriteMany (RWX)}: The volume can be mounted as read-write by many nodes simultaneously. Requires a filesystem or storage backend that supports concurrent writes from multiple hosts (NFS, CephFS, Azure Files, EFS). Many block storage backends do not support RWX.
  \item \textbf{ReadWriteOncePod (RWOP)}: Introduced in Kubernetes 1.22, restricts mounting to a single Pod (not just a single node), providing stricter exclusivity guarantees.
\end{itemize}

A \textbf{PersistentVolumeClaim (PVC)} is a namespaced resource representing a Pod's request for storage. A PVC specifies desired capacity, access mode, and optionally a StorageClass. Kubernetes matches PVCs to PVs through a binding process: a PVC binds to the first available PV that satisfies its requirements (or, with dynamic provisioning, a new PV is created from a StorageClass template). Once bound, the PVC-PV relationship is exclusive and the Pod can mount the claimed storage as a volume.

\section{StorageClasses and Dynamic Provisioning}
Before StorageClasses, PVs were statically provisioned: an administrator created PVs by hand and users created PVCs hoping a matching PV would exist. This was operationally impractical at scale.

\textbf{StorageClasses} provide templates for dynamic provisioning: when a PVC is created referencing a StorageClass, Kubernetes automatically provisions a new PV matching the PVC's requirements using the StorageClass's provisioner. A StorageClass specifies:

\begin{itemize}
  \item • \textbf{Provisioner}: The plugin or CSI driver responsible for creating volumes. Example: ebs.csi.aws.com, disk.csi.azure.com, rook-ceph.rbd.csi.ceph.com.
  \item \textbf{Parameters}: Storage backend-specific configuration (volume type, IOPS, encryption key, etc.).
  \item \textbf{Reclaim policy}: What happens to the PV when the PVC is deleted --- Retain (keep the PV and its data for manual cleanup) or Delete (automatically delete the PV and its underlying storage resource).
  \item \textbf{Volume binding mode}: WaitForFirstConsumer delays PV provisioning until a Pod claiming the PVC is scheduled, allowing the provisioner to create the volume in the same availability zone as the Pod.
  \item \textbf{Volume expansion}: Whether users can resize PVCs after creation.
\end{itemize}

A cluster may have multiple StorageClasses representing different performance tiers (fast SSD, capacity HDD, NFS for shared access), and one StorageClass can be designated as the default, used when a PVC does not specify a class.

\section{StatefulSets and Stable Storage Identity}
For stateful applications --- databases, message queues, distributed systems with node identity requirements --- the \textbf{StatefulSet} workload resource provides the stable identity guarantees that Deployments cannot. Unlike Deployment Pods (which are interchangeable and have ephemeral identities), StatefulSet Pods have:

\begin{itemize}
  \item • \textbf{Stable, unique pod names}: pod-0, pod-1, pod-2, persistent across restarts and rescheduling.
  \item \textbf{Stable network identity}: Headless service DNS entries (pod-0.service.namespace.svc.cluster.local) that persist across pod reschedules.
  \item \textbf{Stable storage}: Each StatefulSet Pod has a dedicated PVC bound to it, created from a \textbf{volumeClaimTemplate} specified in the StatefulSet manifest. When pod-0 is deleted and rescheduled on a different node, it reattaches to the same PVC and PV that it previously used. This is essential for databases (each replica has its own data directory that must follow the replica across reschedules).
\end{itemize}

StatefulSets also control Pod creation and deletion ordering --- Pods are created sequentially (0, 1, 2...) and deleted in reverse order --- and support rolling updates with configurable partition settings, allowing staged rollouts.

\section{The Container Storage Interface (CSI)}
Before CSI, Kubernetes included storage drivers for popular cloud providers and storage systems directly in the Kubernetes codebase ("in-tree" drivers). This meant that adding a new storage system required modifying Kubernetes itself, coordinating with the Kubernetes release cycle, and waiting for a new Kubernetes version to ship. Updating a storage driver to fix a bug required upgrading the entire Kubernetes cluster.

\textbf{Container Storage Interface (CSI)} is a vendor-neutral specification defining a standard RPC interface between container orchestrators (Kubernetes, Mesos, Docker Swarm) and storage drivers. CSI drivers are deployed as separate Pods in the Kubernetes cluster, running alongside the Kubernetes control plane but decoupled from it. Storage vendors can release new driver versions, bug fixes, and features independently of the Kubernetes release cycle.

\section{CSI Driver Architecture}

Figure~\ref{fig:k8s-csi-pvpvc} illustrates the dynamic provisioning flow from a Pod's PVC request through the StorageClass and CSI driver to the underlying storage backend, showing how Kubernetes binds PVCs to dynamically provisioned PersistentVolumes.

\begin{figure}[htbp]
\centering
\begin{tikzpicture}[scale=0.88, every node/.style={transform shape}]
  % Top row: Pod and PVC
  \node[component dark, minimum width=2.8cm, minimum height=1cm] (pod) at (-3,4.5) {Pod};
  \node[component accent, minimum width=3.2cm, minimum height=1cm] (pvc) at (3,4.5) {PVC};
  \node[smalllabel] at (3,3.8) {10Gi, RWO};

  % Arrow Pod -> PVC
  \draw[dataarrow] (pod.east) -- (pvc.west)
    node[midway, above, font=\scriptsize\bfseries] {\textcircled{\scriptsize 1} mount};

  % Middle row: StorageClass and CSI Driver
  \node[component, minimum width=3.2cm, minimum height=1cm] (sc) at (-3,2.2) {StorageClass};
  \node[smalllabel] at (-3,1.5) {trident-gold};
  \node[component dark, minimum width=3.5cm, minimum height=1cm] (csi) at (3,2.2) {CSI Driver (Trident)};

  % Arrow StorageClass -> CSI Driver
  \draw[dataarrow] (sc.east) -- (csi.west)
    node[midway, above, font=\scriptsize\bfseries] {\textcircled{\scriptsize 2} provision};

  % Arrow PVC -> StorageClass (references)
  \draw[dataarrow thin, dashed] (pvc.south) -- (csi.north)
    node[midway, right, font=\scriptsize] {trigger};

  % Bottom row: PV and ONTAP FlexVol
  \node[component accent, minimum width=3.2cm, minimum height=1cm] (pv) at (-3,-0.2) {PersistentVolume};
  \node[component dark, minimum width=3.2cm, minimum height=1cm] (ontap) at (3,-0.2) {ONTAP FlexVol};

  % Arrow CSI -> ONTAP
  \draw[dataarrow] (csi.south) -- (ontap.north)
    node[midway, right, font=\scriptsize\bfseries] {\textcircled{\scriptsize 3} create vol};

  % Arrow PV -> ONTAP
  \draw[dataarrow thin] (pv.east) -- (ontap.west)
    node[midway, above, font=\scriptsize] {backed by};

  % Arrow PVC -> PV (bind)
  \draw[dataarrow] (pvc.south west) -- (pv.north east)
    node[midway, left, font=\scriptsize\bfseries] {\textcircled{\scriptsize 4} bind};

  % Numbered circle helper
  % (using a simple approach with circled numbers in text)
\end{tikzpicture}
\caption{Kubernetes CSI dynamic provisioning flow: (1) a Pod mounts a PVC, (2) the StorageClass triggers the CSI driver to provision storage, (3) the CSI driver creates a volume on the ONTAP backend, and (4) Kubernetes binds the PVC to the resulting PersistentVolume.}
\label{fig:k8s-csi-pvpvc}
\end{figure}

A CSI driver consists of two components:

\textbf{Controller Plugin} runs as a Deployment (typically one or a small number of replicas for HA) in a privileged namespace. It handles cluster-scoped operations: volume creation and deletion, volume snapshots, volume attachment/detachment to nodes, and volume expansion. The Controller Plugin communicates with the storage backend's API (AWS, Azure, Ceph, etc.) to perform these operations. Kubernetes's external-provisioner, external-attacher, external-snapshotter, and external-resizer sidecar containers run alongside the Controller Plugin, watching Kubernetes API objects and calling the appropriate CSI RPC methods.

\textbf{Node Plugin} runs as a DaemonSet, meaning one instance runs on every Kubernetes node. It handles node-scoped operations: mounting a volume to the host's filesystem, staging (bind-mounting), and unmounting volumes when Pods terminate. The node-registrar sidecar registers the Node Plugin with the kubelet on each node.

The kubelet calls the Node Plugin's CSI RPCs to mount and unmount volumes as Pods are scheduled and terminated. The separation between Controller Plugin (cluster-wide volume lifecycle) and Node Plugin (per-node mounting) allows the CSI driver to implement operations that require access to both the storage backend API (creating volumes) and the node's local storage stack (mounting block devices).

\section{Volume Snapshots and Clones via CSI}
CSI extended the storage abstraction to include \textbf{VolumeSnapshots}. A VolumeSnapshot (and the associated VolumeSnapshotClass and VolumeSnapshotContent resources) allows applications to request point-in-time snapshots of PVCs through standard Kubernetes API objects. A snapshot can then be used as the data source for a new PVC, enabling efficient copy-on-write clones of database volumes --- spinning up a new database replica from a snapshot rather than a full data copy.

\textbf{Volume cloning} is a related CSI feature where a new PVC is created with an existing PVC as its data source. If the storage backend supports efficient cloning (COW clone rather than full data copy), this operation can be near-instantaneous regardless of the data volume.

\section{Container-Native Storage Platforms}
Container-native storage platforms are storage systems designed specifically to run inside Kubernetes clusters, using Kubernetes-native operational models (Operators, CRDs, CSI) and leveraging the cluster's own compute resources for storage services.

\section{Portworx (Pure Storage)}
Portworx is a commercial container-native storage platform that runs as a DaemonSet on every Kubernetes node, aggregating local storage into a cluster-wide distributed storage pool. Every node with local drives contributes to the pool; Portworx then provisions PVs from this pool through its CSI driver.

Portworx's key differentiator is its comprehensive data services layer on top of the raw distributed storage:

\textbf{Application-aware snapshots} integrate with Kubernetes application resources. Portworx can create consistent, crash-consistent snapshots of PVCs associated with a specific application (for example, all PVCs in a StatefulSet representing a distributed database) in a coordinated fashion, ensuring that the snapshot captures a consistent application state rather than independent per-volume snapshots taken at slightly different times.

\textbf{Async and sync DR}: Portworx replicates PVC data asynchronously to a secondary Kubernetes cluster for disaster recovery, with the ability to fail over applications (recreating StatefulSets and their PVCs on the DR cluster) through a runbook-driven process. Synchronous replication provides near-zero RPO for critical workloads.

\textbf{Topology awareness}: Portworx understands Kubernetes topology labels (availability zones, regions) and can enforce placement rules ensuring that storage replicas span availability zones, providing resilience to zone failures without relying on cloud provider block storage features.

\textbf{Storage-level encryption} with per-PVC or per-namespace keys, integrated with HashiCorp Vault, AWS KMS, Azure Key Vault, and Google Cloud KMS.

Portworx's commercial positioning targets enterprise organizations running mission-critical stateful workloads on Kubernetes --- databases, AI/ML training jobs, financial applications --- where the cost of data loss or extended recovery is high and the built-in storage capabilities of cloud providers or HCI platforms are insufficient for the required service level.

\section{Rook/Ceph}
Rook is a CNCF (Cloud Native Computing Foundation) graduated project that provides a Kubernetes Operator for deploying and managing Ceph clusters inside Kubernetes. Rook translates Kubernetes-style declarative resource definitions (CRDs like CephCluster, CephBlockPool, CephFilesystem, CephObjectStore) into Ceph cluster configuration and operations.

Deploying Ceph through Rook provides the full Ceph feature set (RADOS block storage via RBD, shared filesystem via CephFS, object storage via RGW) through CSI drivers, with lifecycle management handled by the Rook Operator. The operator monitors Ceph cluster health, handles OSD addition and removal, manages CephFS MDS failover, and upgrades the Ceph version in place.

Rook/Ceph is particularly compelling for organizations running Kubernetes on bare metal or on hardware where cloud-provided storage is unavailable or too expensive. The combination of Ceph's mature distributed storage with Rook's Kubernetes-native management provides a full-featured storage platform that rivals dedicated external storage arrays in capability, at a fraction of the cost when run on commodity hardware.

The operational challenge with Rook/Ceph is that operating a distributed Ceph cluster inside Kubernetes adds complexity: the Kubernetes cluster must be healthy for the storage to be healthy, yet storage failures can impact Kubernetes cluster health, creating potential circular dependency issues. Running Ceph control plane components (monitors, managers) on nodes dedicated to storage rather than on general-purpose nodes is recommended for production deployments.

\section{OpenEBS}
OpenEBS takes a modular approach to container storage, offering multiple storage engines within a single framework, allowing administrators to choose the most appropriate engine for each workload:

\textbf{Jiva}: A simple replication-based engine implemented entirely in user space. Jiva is easy to deploy and understand but has limited performance compared to other engines. It is appropriate for development environments and workloads with modest I/O demands.

\textbf{cStor}: A storage engine backed by OpenZFS's storage pool management. cStor provides ZFS-like features (checksums, snapshots, clone, thin provisioning) within Kubernetes PVs, with replication across nodes. cStor pools can be composed of block devices on each node.

\textbf{Mayastor (now OpenEBS Replicated Engine)}: The highest-performance engine, built using SPDK and NVMe-oF, with Mayastor I/O controllers processing storage requests entirely in user space with direct NVMe hardware access. Mayastor replicates data synchronously across NVMe devices on different nodes over an NVMe-oF fabric (using TCP or RDMA transports). The result is near-native NVMe performance with replication and Kubernetes-native management. Mayastor targets latency-sensitive database workloads where the overhead of other storage engines would be unacceptable.

\textbf{LocalPV (hostpath and device)}: For workloads that manage their own data replication (distributed databases like Cassandra, Kafka, TiDB), OpenEBS provides local persistent volumes that attach a specific disk or filesystem path to a Pod without any additional replication layer. The application handles redundancy; OpenEBS handles the Pod-to-node affinity (ensuring the Pod is always scheduled on the node with its data).

\section{NetApp Trident and Astra}
NetApp's Kubernetes storage integrations are organized under two product lines:

\textbf{NetApp Trident} is a CSI driver that integrates Kubernetes workloads with NetApp's on-premises and cloud storage systems: ONTAP (NFS, iSCSI, NVMe/TCP, NVMe/FC), Element (iSCSI), and cloud-based ONTAP (Azure NetApp Files, Amazon FSx for NetApp ONTAP, Cloud Volumes ONTAP). Trident handles dynamic PVC provisioning, volume snapshots, cloning, and volume expansion. For NFS-backed PVCs, Trident can provision both FlexVol volumes (one Kubernetes volume per ONTAP FlexVol, with strong isolation) and qtrees (multiple Kubernetes volumes sharing a single ONTAP FlexVol, for higher density). For block-backed PVCs, Trident provisions iSCSI or NVMe LUNs.

Trident's value for organizations with existing NetApp infrastructure is continuity: the same ONTAP system, data management policies, and operational workflows that serve traditional workloads can also serve Kubernetes applications, without requiring separate storage infrastructure.

\textbf{NetApp Astra Control} adds a data management layer on top of Trident, providing application-aware backup, restore, cloning, and migration capabilities for Kubernetes applications. Astra Control understands the full application topology (PVCs, ConfigMaps, Secrets, network policies) and can capture and restore it as a complete unit. This is the distinction from simple PVC backups: Astra Control captures the entire application state, enabling migration of a running application (StatefulSet + PVCs + associated resources) from one cluster to another, including cross-cloud migrations.

\section{Stateful Workloads on Kubernetes}
Running stateful workloads on Kubernetes was once considered experimental or even inadvisable. As the ecosystem has matured, specific patterns and tooling have emerged that make databases, message queues, and data analytics engines production-ready on Kubernetes.

\section{Databases on Kubernetes}
Relational databases (PostgreSQL, MySQL, SQL Server) have been successfully deployed on Kubernetes at production scale using StatefulSets with block storage (RWO PVCs), with database replication handled by the database itself. Each replica is a StatefulSet pod with its own PVC containing the replica's data directory. Storage I/O performance depends heavily on the underlying PV implementation: local persistent volumes (hostPath or local device) provide the best performance (near-bare-metal NVMe) but sacrifice Kubernetes's standard scheduling flexibility; cloud block storage (EBS, Azure Premium SSD) provides acceptable latency for many database workloads with the benefit of easy multi-AZ operation.

\textbf{Kubernetes Operators} have been critical to making databases on Kubernetes operationally viable. Operators encode the operational knowledge for specific databases: the PostgreSQL operator knows how to initialize replica clusters, handle primary failures and promote standbys, manage connection pooling, and coordinate backups. Without operators, all of this would require manual intervention for every database lifecycle event.

\textbf{CloudNativePG} (for PostgreSQL), \textbf{Percona Operator} (for MySQL and MongoDB), \textbf{Strimzi} (for Kafka), and \textbf{KubeDB} are examples of mature operators that have pushed database-on-Kubernetes into production at substantial scale.

\section{AI/ML Training Workloads and Storage}
AI/ML training workloads place some of the most demanding storage requirements of any workload class on Kubernetes. A large language model training job may involve:

\begin{itemize}
  \item • Hundreds or thousands of GPUs across many nodes, all requiring simultaneous access to a shared training dataset
  \item Reading data at rates that can saturate multi-terabit-per-second storage systems
  \item Writing checkpoints (periodically saving model state to resume from if a job fails) at intervals of potentially hundreds of gigabytes every few minutes
\end{itemize}

For the training data read path, the storage system must deliver very high aggregate bandwidth to many readers simultaneously --- this is the use case for parallel filesystems (Lustre, GPFS, Weka) accessed via Kubernetes PVCs, or for cloud object storage (S3) accessed directly by the training framework rather than through the Kubernetes storage layer. AWS's \textbf{FSx for Lustre}, with S3 as the data lake backing, is a common choice for managed high-bandwidth training data access on AWS.

For checkpointing, the storage system must be able to absorb a burst of large sequential writes at high bandwidth. A NFS or SMB share backed by high-performance NAS, or an RBD volume backed by a Ceph cluster with high-endurance SSDs, is typically appropriate.

\textbf{GPUDirect Storage} (NVIDIA's DMA-based storage access protocol, avoiding CPU and system memory for storage transfers) is increasingly relevant in large AI training clusters. Storage backends that support GPUDirect --- WekaFS, IBM Spectrum Scale, and some NVMe-oF configurations --- allow GPU memory to receive training data directly from storage without CPU involvement, reducing I/O latency and freeing CPU cycles for training computation.

\section{Multi-Cloud Storage Abstraction and Data Mobility}
As organizations spread workloads across multiple cloud providers and between cloud and on-premises infrastructure, the question of data portability --- moving data between environments without excessive cost, complexity, or delay --- becomes critical.

\section{The Multi-Cloud Storage Challenge}
Cloud storage services are deeply integrated with their respective cloud ecosystems: EBS volumes attach only to EC2 instances, Azure Managed Disks attach only to Azure VMs, and moving data between clouds requires first exporting it to object storage, transferring it across cloud networks (at data egress costs), and re-importing it. This creates a practical lock-in that is less about contractual constraints and more about engineering friction and cost.

Object storage has become the de facto multi-cloud data exchange format because S3 API compatibility is near-universal: AWS S3, Azure Blob Storage (with S3-compatibility layer), Google Cloud Storage, MinIO, and dozens of on-premises object stores all speak variants of the S3 API. Datasets stored as objects can be accessed by applications in any cloud with an S3-compatible client, without conversion. This is why S3-compatible object storage has become the standard landing zone for data lakes, model repositories, and backup archives in multi-cloud architectures.

\section{Kubernetes Storage Abstraction for Multi-Cloud}
For Kubernetes workloads spanning multiple clouds, the CSI model provides a partial abstraction: a PVC backed by an EBS volume on AWS and a PVC backed by an Azure Managed Disk on Azure are both PVCs from the application's perspective. If the application's YAML manifests avoid cloud-provider-specific annotations and rely on StorageClass names (which can be configured to point to different backends in different clusters), the application manifests can be relatively portable.

The practical challenge is data state: the application may be portable, but its data is stored in a cloud-specific PV. Migrating a stateful application between clouds requires migrating its data, which can be accomplished through:

\begin{itemize}
  \item • \textbf{Storage vendor replication}: Portworx's async DR replicates PVC data between clusters in different clouds. NetApp Astra Control replicates application state (including PVCs) across clusters.
  \item \textbf{Object storage intermediary}: Application-level backup to S3-compatible object storage, restore in the target cloud. This is simple but requires application downtime during the restore.
  \item \textbf{Database-level replication}: For database workloads, using the database's own replication mechanism (PostgreSQL streaming replication, MySQL GTID replication, Kafka MirrorMaker) to replicate data to a standby in the target cloud, then promoting the standby during migration.
\end{itemize}

\section{Data Fabric and Abstraction Platforms}
Several vendors offer storage abstraction platforms that present a unified namespace across multiple clouds and on-premises environments:

\textbf{NetApp Cloud Volumes ONTAP} runs ONTAP software in cloud VMs (EC2, Azure VMs, GCE), providing ONTAP-native NFS and iSCSI to cloud workloads. SnapMirror replication can then transfer data between on-premises ONTAP systems and cloud-based ONTAP systems, and between ONTAP instances in different clouds.

\textbf{Qumulo's Scale Anywhere} architecture deploys Qumulo software both on-premises and on cloud instances (using cloud block storage as the backing), presenting a single namespace that spans both environments. A file created on-premises can appear in the cloud namespace within minutes, enabling hybrid workflows without a separate data movement step.

\textbf{VAST Data's global flash architecture} operates within a single cluster but can be extended through VAST's global namespace federation, allowing multiple VAST clusters (potentially in different sites or clouds) to share a directory namespace where files physically located in one cluster appear in the namespace of another.

For the broadest multi-cloud data portability, the pragmatic solution remains standardization on S3-compatible object storage for unstructured data, combined with application-level data replication for stateful services, and reliance on managed database services (RDS, Azure Database, Cloud SQL) that abstract the storage layer from the application entirely.

\section{Security \& Governance}
Kubernetes storage security requires attention at multiple layers: the storage data itself, the control plane that provisions and manages storage, the secrets used to authenticate to storage backends, and the encryption of data in transit and at rest.

\section{CSI Encryption}
CSI drivers can implement encryption at two levels: storage-side encryption (the storage backend encrypts data before persisting it) and in-transit encryption (data is encrypted between the Kubernetes node and the storage backend).

\textbf{Storage-side encryption} for CSI volumes is typically implemented through StorageClass parameters that instruct the storage backend to encrypt the provisioned volume. For AWS EBS, the StorageClass parameter encrypted: "true" and an optional kmsKeyId enable EBS encryption using AWS KMS. For Azure Managed Disks, disk encryption sets (DES) that reference Azure Key Vault keys are specified in the StorageClass. For Ceph RBD via Rook, per-PVC encryption can be enabled with the key stored in a Kubernetes Secret or in an external KMS (HashiCorp Vault via the Vault KMS plugin).

\textbf{In-transit encryption} between the Kubernetes node and the storage backend depends on the storage protocol. For NFS-backed PVs (EFS, FSx, Azure Files, CephFS), NFSv4 with Kerberos (NFS sec=krb5p) encrypts data in transit. For block storage accessed over iSCSI, enabling IPsec between initiators and targets (typically at the cluster or node level rather than per-volume) encrypts replication traffic. For CSI drivers using proprietary protocols (Portworx, NetApp Trident iSCSI), vendor-specific TLS configuration applies.

\textbf{LUKS encryption at the node level} is available through some CSI drivers (including Rook/Ceph and OpenEBS Mayastor): the driver formats the block device with LUKS encryption at the node before mounting it, and the encryption key is retrieved from an external KMS. This approach encrypts data even if the underlying storage backend does not support native encryption, and places encryption key management fully under the Kubernetes cluster's control.

\section{Kubernetes RBAC for Storage}
Kubernetes RBAC (Role-Based Access Control) governs which users and service accounts can create, modify, or delete PVs, PVCs, StorageClasses, and VolumeSnapshots. Proper RBAC configuration for storage resources is frequently overlooked in favor of application-level RBAC.

Key RBAC considerations for storage:

\textbf{PVC creation}: In most clusters, any authenticated user in a namespace can create PVCs. For multi-tenant clusters, limiting PVC creation (through namespace-specific RBAC, admission controllers, or quota enforcement) prevents tenants from requesting excessive storage or consuming StorageClasses they should not access.

\textbf{StorageClass access}: By default, all StorageClasses are cluster-global and accessible to all namespaces. For multi-tenant environments, Kubernetes does not natively restrict which StorageClasses a namespace can use (unlike ResourceQuotas which can limit storage capacity). This gap is typically addressed through admission webhooks (Open Policy Agent/Gatekeeper, Kyverno) that enforce namespace-to-StorageClass binding rules.

\textbf{VolumeSnapshot access}: VolumeSnapshots are namespace-scoped, and RBAC should limit who can create or restore from snapshots in each namespace. A developer who can create snapshots of production database PVCs represents a data exfiltration risk --- the snapshot can be restored into a development namespace and the sensitive data accessed without production-level controls.

\textbf{CSI driver service accounts}: The Controller Plugin and Node Plugin service accounts require specific Kubernetes permissions to manage PVs and interact with node resources. These permissions should be scoped to exactly what the CSI driver requires, using dedicated ClusterRoles rather than granting cluster-admin.

\section{Secrets Management for Storage Credentials}
CSI drivers authenticate to storage backends using credentials --- AWS access keys, Azure service principal credentials, Ceph admin keys, NetApp SVM credentials. These credentials are typically stored in Kubernetes Secrets in the CSI driver's namespace. Kubernetes Secrets are base64-encoded (not encrypted) in etcd by default, meaning that any actor with access to etcd --- or with kubectl get secret permission in the driver's namespace --- can read storage backend credentials.

Mitigating this risk requires:

\textbf{etcd encryption at rest}: Kubernetes supports encrypting Secret resources in etcd using AES-CBC or AES-GCM encryption providers. The encryption key is held by the kube-apiserver, which decrypts Secrets on read. Alternatively, an external KMS (AWS KMS, Azure Key Vault, GCP Cloud KMS, HashiCorp Vault) can be integrated through the Kubernetes KMS plugin framework, where etcd stores envelope-encrypted Secrets (the Secret's data is encrypted by a data encryption key, which itself is encrypted by the external KMS). This ensures that the encryption key never resides on the Kubernetes nodes themselves.

\textbf{External secrets operators}: The External Secrets Operator (ESO) synchronizes secrets from external secrets managers (HashiCorp Vault, AWS Secrets Manager, Azure Key Vault, GCP Secret Manager) into Kubernetes Secrets, providing a lifecycle management layer that includes automatic rotation. The CSI driver's credentials are rotated in the external secrets manager, and the operator automatically updates the Kubernetes Secret, avoiding long-lived credentials in etcd.

\textbf{Workload identity}: For cloud-hosted Kubernetes (EKS, AKS, GKE), the most secure credential pattern for cloud-native CSI drivers is workload identity --- assigning an IAM role or Azure managed identity to the CSI driver's service account, allowing the driver to authenticate to cloud storage services using cryptographic pod identity rather than static credentials. AWS EKS Pod Identity, Azure Workload Identity, and GKE Workload Identity all implement this pattern, eliminating the need to store cloud credentials as Kubernetes Secrets entirely.

\section{Cloud KMS Integration}
All major cloud block and file storage services support integration with their respective KMS offerings for server-side encryption key management:

\textbf{AWS KMS} manages customer master keys (CMKs) used to encrypt EBS volumes, EFS filesystems, and S3 buckets. EBS volumes can use AWS-managed keys (aws/ebs service key, managed automatically) or customer-managed keys (CMKs) where the customer controls key policy, rotation schedule, and key deletion. KMS key policies govern which IAM principals (EC2 instances, Lambda functions, EKS pods) can use the key to encrypt or decrypt data. Using CMKs provides audit logging of every key usage through AWS CloudTrail, enabling visibility into who accessed what data and when.

\textbf{Azure Key Vault} serves an equivalent role for Azure Managed Disks (through Disk Encryption Sets), Azure Files (through customer-managed keys in storage accounts), and other Azure storage services. Azure Key Vault supports both software-protected keys and HSM-protected keys (using FIPS 140-2 Level 3 certified HSMs through Azure Dedicated HSM or Key Vault Managed HSM). The ability to use HSM-protected keys for storage encryption at cloud scale is a significant capability for organizations in regulated industries.

\textbf{Google Cloud KMS} manages CMEK (Customer-Managed Encryption Keys) for Persistent Disks, Filestore, and GCS. GCP's Cloud EKM (External Key Manager) integration allows customers to use keys held in external HSMs (Thales, Fortanix, Futurex) while encrypting GCP storage resources, implementing a HYOK (Hold Your Own Key) model where encryption keys never reside in Google infrastructure.

\textbf{HashiCorp Vault} is widely used as a cloud-agnostic KMS integration point, particularly in multi-cloud environments where per-cloud KMS creates management fragmentation. Vault's transit secrets engine provides encryption-as-a-service and can generate data encryption keys for storage systems without ever exposing the master key to the consuming application. The Vault Agent Injector pattern injects Vault-sourced secrets into Pods as files or environment variables without storing them in Kubernetes etcd.

For any Kubernetes cluster handling sensitive data, enabling etcd encryption at rest with external KMS integration, using workload identity for cloud storage authentication, and enforcing PVC-level encryption through StorageClass parameters should be considered baseline security controls rather than optional hardening.

